\documentclass{article}
\usepackage[T1]{fontenc}
\usepackage{lmodern}
\usepackage{graphicx}
\usepackage{amsmath,amssymb}
\usepackage[a4paper,margin=2cm]{geometry}
\usepackage[absolute,overlay]{textpos}
\setlength{\TPHorizModule}{1mm}
\setlength{\TPVertModule}{1mm}
\usepackage[indonesian]{babel}
\usepackage{hyperref}
\setlength{\emergencystretch}{2em}
% unit: Halaman 21

\begin{document}

% === Halaman 21 (llm) ===

\section*{Implementasi Knapsack}

\begin{itemize}
    \item Ukuran cipherteks yang dihasilkan lebih besar daripada plainteksnya, karena enkripsi dapat menghasilkan kriptogram yang nilai desimalnya lebih besar daripada nilai desimal blok plainteks yang dienkripsikan.
    \item Untuk menambah kekuatan algoritma \textit{knapsack}, kunci publik maupun kunci privat seharusnya paling sedikit 250 elemen, nilai setiap elemen antara 200 sampai 400 bit panjangnya, nilai modulus antara 100 sampai 200 bit.
    \item Dengan nilai-nilai \textit{knapsack} sepanjang itu, dibutuhkan $10^{46}$ tahun untuk menemukan kunci secara \textit{brute force}, dengan asumsi satu juta percobaan setiap detik.
\end{itemize}

\end{document}
